% 附录 F 部分习题答案与提示（Answers and hints to selected problems）
\chapter{部分习题答案与提示}

\exer{1.1} $9.27\times10^{-24}$ A\,m$^{2}$；8.40 GHz；34.7 $\mu$eV；8.40 GHz。这里我们只考虑
自旋角动量。

\exer{1.4} 这是标量算符 $(\hat{\mathbf{S}}\cdot\mathbf{X})$ 与矢量算符
$(\hat{\mathbf{S}})$ 的对易子。算出它的一个分量，即可求值
\begin{align*}
[\hat{\mathbf{S}}\cdot\mathbf{X}, \hat{\mathbf{S}}\cdot\mathbf{Y}] &= \frac{1}{4}\left((\boldsymbol{\sigma}\cdot\mathbf{X})(\boldsymbol{\sigma}\cdot\mathbf{Y}) - (\boldsymbol{\sigma}\cdot\mathbf{y}) (\boldsymbol{\sigma}\cdot\mathbf{X})\right)\\
&= \frac{\mathrm{i}}{4}\left(\boldsymbol{\sigma}\cdot(\mathbf{X}\times\mathbf{Y}) - \boldsymbol{\sigma}\cdot(\mathbf{Y}\times\mathbf{X})\right)\\
&= \frac{\mathrm{i}}{2}\boldsymbol{\sigma}\cdot(\mathbf{X}\times\mathbf{Y})\\
&= \mathrm{i}\mathbf{S}\times\mathbf{X}\cdot\mathbf{Y},\tag{F.1}
\end{align*}
结果得证。\footnote{首行 $(\boldsymbol{\sigma}\cdot\mathbf{y})$ 原书如此，应为
$(\boldsymbol{\sigma}\cdot\mathbf{Y})$——疑为原书排印之误。末行 $\mathbf{S}$ 上的
帽子原书即未加，亦照录。——译注}

\exer{1.5} 用式 1.59 可以证明
\begin{equation*}
\hat{S}_z(\hat{S}_{\pm}|S, S_z\rangle) = (\hat{S}_{\pm}\hat{S}_z \pm \hat{S}_{\pm})|S, S_z\rangle = (S_z\pm1)(\hat{S}_{\pm}|S, S_z\rangle),\tag{F.2}
\end{equation*}
即自旋的 z 分量被升高或降低。式 1.58 与 1.61 可以写作
\begin{equation*}
\hat{S}_{+}\hat{S}_{-} - \hat{S}_{-}\hat{S}_{+} = 2\hat{S}_z\tag{F.3}
\end{equation*}
\begin{equation*}
\hat{S}_{+}\hat{S}_{-} + \hat{S}_{-}\hat{S}_{+} = 2(\hat{S}_x^{2}+\hat{S}_y^{2}) = \hat{\mathbf{S}}^{2}-\hat{S}_z^{2},\tag{F.4}
\end{equation*}
相加或相减得
\begin{align*}
\hat{S}_{+}\hat{S}_{-} &= \hat{\mathbf{S}}^{2}-\hat{S}_z^{2}+\hat{S}_z\tag{F.5}\\
\hat{S}_{-}\hat{S}_{+} &= \hat{\mathbf{S}}^{2}-\hat{S}_z^{2}-\hat{S}_z.\tag{F.6}
\end{align*}
故
\begin{align*}
\langle S, S_z|\hat{S}_{-}\hat{S}_{+}|S, S_z\rangle &= S(S+1)-S_z^{2}-S_z\\
&= S(S+1)-S_z(S_z+1)\tag{F.7}
\end{align*}
及
\begin{align*}
\langle S, S_z|\hat{S}_{+}\hat{S}_{-}|S, S_z\rangle &= S(S+1)-S_z^{2}+S_z\\
&= S(S+1)-S_z(S_z-1)\tag{F.8}
\end{align*}
所需的归一化于是得证。

\exer{1.6} 不失一般性，取 $\mathbf{B}=(0,0,B)$，则 $\mathbf{B}\times\mathbf{r}=B(-y,-x,0)$，
$\frac{1}{2}\nabla\times(\mathbf{B}\times\mathbf{r})=(0,0,B)$。其他选取也可行，例如
$\mathbf{A}=-By(1,0,0)$ 或 $\mathbf{A}=Bx(0,1,0)$。事实上
$\mathbf{A}=B((\alpha-1)y,\alpha x,0)$ 都可行，$\alpha$ 为任意实数。

\exer{1.7} 提示：计算 $[\boldsymbol{\sigma}\cdot(\mathbf{p}+e\mathbf{A})]^{2}/2m_{\mathrm{e}}$。

\exer{1.8} 不失一般性，新轴可取在 $xz$ 平面内、与 $z$ 轴成 $\theta$ 角。算符为
\begin{equation*}
\hat{S}_{\theta} = \cos\theta\hat{S}_z + \sin\theta\hat{S}_x.\tag{F.9}
\end{equation*}
本征值为 1 的本征矢为
\begin{equation*}
(2(1-\cos\theta))^{-1/2}\begin{pmatrix}\sin\theta\\1-\cos\theta\end{pmatrix},\tag{F.10}
\end{equation*}
故所求概率为
\begin{equation*}
\frac{\sin^{2}\theta}{2(1-\cos\theta)} = \cos^{2}\frac{\theta}{2}.\tag{F.11}
\end{equation*}

\exer{1.10}
\begin{equation*}
H_{\mathrm{d}}/H_{\mathrm{a}}:\ \text{(a)}\ 3\times10^{-5}\quad \text{(b)}\ -8.8\times10^{-4}
\end{equation*}
\begin{equation*}
B_{\mathrm{d}}/B_{\mathrm{a}}:\ \text{(a)}\ -6\times10^{-5}\quad \text{(b)}\ 1.76\times10^{-3}
\end{equation*}

\exer{1.12} 提示：$\sigma_m^{2}=I$，故
\begin{equation*}
\sigma^{n} = \begin{cases} I & n \text{ 为偶}\\ \sigma_m & n \text{ 为奇}. \end{cases}\tag{F.12}
\end{equation*}
于是
\begin{align*}
\mathrm{e}^{\mathrm{i}\alpha\sigma_m} &= I\left(1-\frac{\alpha^{2}}{2!}+\frac{\alpha^{4}}{4!}-\cdots\right)\\
&\quad + \mathrm{i}\sigma_m\left(\alpha-\frac{\alpha^{3}}{3!}+\frac{\alpha^{5}}{5!}-\cdots\right)\\
&= I\cos\alpha + \mathrm{i}\sigma_m\sin\alpha.\tag{F.13}
\end{align*}

\exer{1.14} （a）力矩 $=I\ddot{\theta}=-pE\sin\theta$，故
$\ddot{\theta}+pE\sin\theta/I=0$；$\theta\ll1$ 时给出简谐运动，$\omega^{2}=pE/I$。
（b）哈密顿量为 $H=\hat{L}^{2}/2I-pE\cos\theta$。此计算中需求出若干对易子：
\begin{equation*}
[\hat{\theta}, \hat{L}] = \mathrm{i}\hbar\tag{F.14}
\end{equation*}
\begin{equation*}
[\hat{\theta}, \hat{L}^{2}] = 2\mathrm{i}\hbar\hat{L}\tag{F.15}
\end{equation*}
\begin{equation*}
[\hat{L}, \cos\theta] = \mathrm{i}\hbar\sin\theta,\tag{F.16}
\end{equation*}
结果随之而来。电偶极子不同于磁矩之处在于它不与角动量相联系（电偶极子只是空间
分开的两个相反电荷；磁偶极子则是相互绕行的两个相反电荷，因而与角动量相联系）。
注意哈密顿量不含耗散，故电偶极子永远振荡、磁偶极子永远进动。只有允许它们与热浴
交换能量（耗散过程）而改变自身能量时，两种偶极子才会沿外场排列。

\exer{2.1} $-1.65\times10^{-15}$；$7.8\times10^{-12}$。

\exer{2.2} 鸭子主要是水，质量 2--3 kg。一枚铁屑约 $0.1\,\mathrm{mm^{3}}=10^{-10}\,\mathrm{m^{3}}$
（估计），每个铁原子携带 $2.2\mu_{\mathrm{B}}$。铁的密度为 $7873\,\mathrm{kg\,m^{-3}}$、
相对原子质量 55.847 g。由此可以算出所需磁场为 1 mT（地磁场约 0.05 mT）。一头牛
约 400 kg，设它也主要由水组成，则质量（从而体积）约大 200 倍，产生同样磁化强度
所需磁场小得多（$\sim2\ \mu$T）。当然我们忽略了血红蛋白的效应。

\exer{2.3} $2.74\times10^{-5}\ \mathrm{J\,T^{-1}}$。

\exer{2.4} 提示：用 $Z=2\cosh(\mu_{\mathrm{B}}B/k_{\mathrm{B}}T)$ 与
$F=-nk_{\mathrm{B}}T\log Z$。则 $E=-n(\mathrm{d}\log Z/\mathrm{d}\beta)$，
$C=\partial E/\partial T$。熵可以用 $S=-(\partial F/\partial T)_B$ 或
$S=(E-F)/T$ 得到。

\exer{2.5} 提示：一个有用的中间结果是
\begin{equation*}
\frac{\partial\log Z_J}{\partial y} = \mathrm{B}_J(y).\tag{F.17}
\end{equation*}

\exer{2.7} （a）${}^{5}\mathrm{I}_8$（b）${}^{4}\mathrm{I}_{15/2}$
（c）${}^{3}\mathrm{H}_6$（d）${}^{1}\mathrm{S}_0$。

\exer{2.8} 提示：（1）你需要计算含 $\mathbf{p}\cdot\mathbf{A}+\mathbf{A}\cdot\mathbf{p}$
的项。它正比于
\begin{equation*}
\mathbf{p}\cdot\boldsymbol{\mu}\times\mathbf{r} + \boldsymbol{\mu}\times\mathbf{r}\cdot\mathbf{p} = 2\mathbf{r}\times\mathbf{p}\cdot\boldsymbol{\mu},\tag{F.18}
\end{equation*}
最后的等号成立是因为 $\boldsymbol{\mu}$ 不依赖位置、与 r 和 p 对易；而 r 与 p
可以互易是因为它们出现在标量三重积中——你从不把 r 与 p 的同一分量相乘。

（2）$\nabla^{2}(\frac{1}{r})=\nabla\cdot\nabla(\frac{1}{r})=4\pi\delta(\mathbf{r})$。
一个类似的结果是 $(\mathbf{S}\cdot\nabla)\nabla(\frac{1}{r})=\mathbf{S}\frac{4\pi}{3}\delta(\mathbf{r})$，
因子 3 是因为 $\mathbf{S}$ 从奇点的三个分量中挑出一个。

\exer{2.9} $1.89\times10^{-4}\ \mathrm{emu\,mol^{-1}}$；$9.68\times10^{-7}\ \mathrm{emu\,g^{-1}}$。

\exer{2.10} 固定 T 时，
\begin{equation*}
T\delta S = -\mathbf{B}\cdot\delta\mathbf{M} - \mathbf{M}\cdot\delta\mathbf{B} + \frac{\partial(nk_{\mathrm{B}}T\log Z)}{\partial\mathbf{B}}\cdot\delta\mathbf{B},\tag{F.19}
\end{equation*}
用 $M=nk_{\mathrm{B}}T(\partial\log Z/\partial B)$ 即得结果。

\exer{3.1} $\mathrm{Sc^{2+}}$ 有一个 3d 电子（$l=2$、$s=\frac{1}{2}$），有 5 个由
$l_z=-2,-1,0,1,2$ 刻画的轨道态。能级为 $0, A, 4A$（计自旋后简并度分别为 2、4、4）；
$A>0$ 时基态二重简并（$E=0$），$A<0$ 时四重简并（$E=4A$）。用记号 $\{l_z, s_z\}$
表示态。$A>0$ 时基态能级为 $\{0,\frac{1}{2}\}$ 与 $\{0,-\frac{1}{2}\}$，不被自旋--
轨道相互作用分裂（因 $\lambda l_zs_z=0$）。任意方向的磁场都会分裂这些能级，在
$\mu_{\mathrm{B}}B\ll k_{\mathrm{B}}T\ll A$ 时给出类居里磁化率。这一条件保证不布居
激发态（$k_{\mathrm{B}}T\ll A$）且留在布里渊函数的线性段（$\mu_{\mathrm{B}}B\ll k_{\mathrm{B}}T$）。

$A<0$ 时基态能级为 $|2,\frac{1}{2}\rangle$、$|2,-\frac{1}{2}\rangle$、
$|-2,\frac{1}{2}\rangle$ 与 $|-2,-\frac{1}{2}\rangle$；自旋--轨道相互作用使
$|2,-\frac{1}{2}\rangle$ 与 $|-2,\frac{1}{2}\rangle$ 能量降低 $\lambda$，
$|2,\frac{1}{2}\rangle$ 与 $|-2,-\frac{1}{2}\rangle$ 升高 $\lambda$。此时只有平行
于 z 的磁场才能分裂较低的能级，故 $\mu_{\mathrm{B}}B\ll k_{\mathrm{B}}T\ll\lambda$
时磁化率是类居里的。磁场垂直于 z 时磁化率不依赖温度。

\exer{3.2} 用
\begin{align*}
(1+x)^{-1/2} &= 1-\frac{x}{2}+\frac{(-\frac{1}{2})(-\frac{3}{2})}{2!}x^{2}+\frac{(-\frac{1}{2})(-\frac{3}{2})(-\frac{5}{2})}{3!}x^{3}\\
&\quad + \frac{(-\frac{1}{2})(-\frac{3}{2})(-\frac{5}{2})(-\frac{7}{2})}{4!}x^{4}+\cdots\\
&= 1-\frac{x}{2}+\frac{3x^{2}}{8}-\frac{5x^{3}}{16}+\frac{35x^{4}}{128}+\cdots.\tag{F.20}
\end{align*}

\exer{3.3} 把 V 化为球坐标：
\begin{equation*}
V = \frac{Dr^{4}}{8}\left[\sin^{4}\theta\left(\mathrm{e}^{4\mathrm{i}\phi}+\mathrm{e}^{-4\mathrm{i}\phi}+6\right) + 8\cos^{2}\theta - \frac{24}{5}\right].\tag{F.21}
\end{equation*}
对矩阵元，除 $V_{2,-2}$ 与 $V_{-2,2}$ 外所有非对角元为零。非零矩阵元的计算直截了当
但繁琐。举一例：
\begin{align*}
V_{1,1} &= \int_{0}^{\infty}r^{6}R^{2}(r)\,\mathrm{d}r\,\frac{D}{8}\int_{0}^{\pi}\left(6\sin^{2}\theta+8\cos^{4}\theta-\frac{24}{5}\right)\\
&\quad\times4\sin^{3}\theta\cos^{2}\theta\,\mathrm{d}\theta\int_{0}^{2\pi}\mathrm{d}\phi\\
&= -4A.\tag{F.22}
\end{align*}
本征值与本征矢为
\begin{center}
\begin{tabular}{ll}
本征值 & 本征矢\\
\midrule
$A+5A\left(1+\frac{4\mu_{\mathrm{B}}^{2}B^{2}}{25A^{2}}\right)^{1/2}$ & $\dfrac{(|2\rangle+x_{+}|-2\rangle)}{(1+x_{+}^{2})^{1/2}}$\\[10pt]
$-4A+\mu_{\mathrm{B}}B$ & $|1\rangle$\\
$6A$ & $|0\rangle$\\
$-4A-\mu_{\mathrm{B}}B$ & $|-1\rangle$\\
$A-5A\left(1+\frac{4\mu_{\mathrm{B}}^{2}B^{2}}{25A^{2}}\right)^{1/2}$ & $\dfrac{(|2\rangle+x_{-}|-2\rangle)}{(1+x_{-}^{2})^{1/2}}$\\
\end{tabular}
\end{center}
其中
\begin{equation*}
x_{\pm} = \pm\left(1+\left(\frac{2\mu_{\mathrm{B}}B}{5A}\right)^{2}\right)^{1/2} - \frac{2\mu_{\mathrm{B}}B}{5A}.\tag{F.23}
\end{equation*}
低场极限下本征值与本征矢为
\begin{center}
\begin{tabular}{ll}
本征值 & 本征矢\\
\midrule
$6A+\dfrac{2\mu_{\mathrm{B}}^{2}B^{2}}{5A}$ & $\dfrac{(|2\rangle+|-2\rangle)}{\sqrt{2}}$\\[10pt]
$-4A+\mu_{\mathrm{B}}B$ & $|1\rangle$\\
$6A$ & $|0\rangle$\\
$-4A-\mu_{\mathrm{B}}B$ & $|-1\rangle$\\
$-4A-\dfrac{2\mu_{\mathrm{B}}^{2}B^{2}}{5A}$ & $\dfrac{(|2\rangle-|-2\rangle)}{\sqrt{2}}$\\
\end{tabular}
\end{center}
故若 $k_{\mathrm{B}}T\ll A$，只布居最低能级：$A>0$ 时是三重态、$A<0$ 时是二重态。
只考虑这些最低能级，配分函数可写为
\begin{equation*}
Z = \begin{cases} 1+e^{-\mu_{\mathrm{B}}B/k_{\mathrm{B}}T}+e^{\mu_{\mathrm{B}}B/k_{\mathrm{B}}T} & A>0\\ 1+e^{-2\mu_{\mathrm{B}}^{2}B^{2}/5Ak_{\mathrm{B}}T} & A<0, \end{cases}\tag{F.24}
\end{equation*}
再用 $F=-Nk_{\mathrm{B}}T\log Z$ 与 $M=-\partial F/\partial B$ 即得结果。

\exer{3.4} $I=\frac{3}{2}$、$A=2\times10^{-6}$ eV。关于核磁矩说不了太多。声子引起
的谱线展宽在高温很大。分辨精细结构常需要低温。

\exer{3.5} 测得 $g_JJ=7$（$\mathrm{Gd^{3+}}$）、5（$\mathrm{Fe^{3+}}$）与 3
（$\mathrm{Cr^{3+}}$）。洪特规则的预言对三者都符合，唯独 $\mathrm{Cr^{3+}}$
例外——按洪特规则应预期 $J=\frac{3}{2}$、$g_J=\frac{2}{5}$。这是因为发生轨道
淬灭、$g_J=2$。（$\mathrm{Fe^{3+}}$ 也发生轨道淬灭，但察觉不到——洪特规则本来
就给出 $L=0$。）

\exer{3.6} 60 MHz 对应的磁场为
\begin{equation*}
\begin{array}{ll}
{}^{1}\mathrm{H} & 1.41\ \mathrm{T}\\
{}^{2}\mathrm{H} & 9.18\ \mathrm{T}\\
{}^{13}\mathrm{C} & 5.60\ \mathrm{T}\\
{}^{19}\mathrm{F} & 1.50\ \mathrm{T}
\end{array}
\end{equation*}

\exer{3.7} 0.322 T。

\exer{3.8} 先导出 $S=1$ 粒子的 $\hat{S}_z$ 与 $\hat{S}_x$ 矩阵表示：
\begin{equation*}
\hat{S}_z = \begin{pmatrix}1 & 0 & 0\\ 0 & 0 & 0\\ 0 & 0 & -1\end{pmatrix}\tag{F.25}
\end{equation*}
\begin{equation*}
\hat{S}_x = \frac{1}{\sqrt{2}}\begin{pmatrix}0 & 1 & 0\\ 1 & 0 & 1\\ 0 & 1 & 0\end{pmatrix}.\tag{F.26}
\end{equation*}
于是哈密顿量为
\begin{equation*}
\hat{\mathcal{H}} = \begin{pmatrix} D+g_{\parallel}\mu_{\mathrm{B}}B\cos\theta & \frac{g_{\perp}\mu_{\mathrm{B}}B\sin\theta}{\sqrt{2}} & 0\\ \frac{g_{\perp}\mu_{\mathrm{B}}B\sin\theta}{\sqrt{2}} & 0 & \frac{g_{\perp}\mu_{\mathrm{B}}B\sin\theta}{\sqrt{2}}\\ 0 & \frac{g_{\perp}\mu_{\mathrm{B}}B\sin\theta}{\sqrt{2}} & D-g_{\parallel}\mu_{\mathrm{B}}B\cos\theta \end{pmatrix}\tag{F.27}
\end{equation*}
求本征值即得所需结果。$\theta=0$ 时本征值简化为
\begin{equation*}
E = 0, D\pm g_{\parallel}\mu_{\mathrm{B}}B.\tag{F.28}
\end{equation*}
$\theta=\pi/2$ 时本征值简化为
\begin{equation*}
E = D, \frac{D}{2}\pm\sqrt{\left(\frac{D}{2}\right)^{2}+(g_{\perp}\mu_{\mathrm{B}}B)^{2}}.\tag{F.29}
\end{equation*}

\exer{3.9} $82.4\,\mathrm{m\,s^{-1}}$；$7.7\times10^{-19}\,\mathrm{m\,s^{-1}}$；
$9.4\times10^{14}\,\mathrm{Hz}$；$9.1\times10^{-9}\,\mathrm{Hz}$；
$2.3\times10^{7}\,\mathrm{Hz}$；$0.1\,\mu\mathrm{eV}$；$7.7\times10^{-4}\,\mathrm{cm^{-1}}$。

\exer{3.10} 若你能轻松测量的最低频率周期为 20 $\mu$s，即频率 50 kHz，对应约
$4\times10^{-4}$ T 的场。活过 20 $\mu$s 的 $\mu$ 子比例 $e^{-20/2.2}=1.1\times10^{-4}$，
即 $10^{7}$ 个注入 $\mu$ 子中约一千出头。

这一结果来自平均寿命 $\tau_{\mu}$ 的 $\mu$ 子指数衰减。活得比 T 久的比分为
\begin{equation*}
\frac{\int_T^{\infty}\mathrm{e}^{-t/\tau_{\mu}}\,\mathrm{d}t}{\int_0^{\infty}\mathrm{e}^{-t/\tau_{\mu}}\,\mathrm{d}t} = \mathrm{e}^{-T/\tau_{\mu}}.\tag{F.30}
\end{equation*}
若脉冲宽度为 50 ns，当自旋进动周期为 100 ns 时会发生相消干涉，对应频率 10 MHz，
给出场上限 $\sim$0.07 T。

若 $B=0.4$ T，$f=54.2$ MHz。

\exer{3.11} 对 $\mathrm{Fe^{3+}}$，$J=\frac{5}{2}$、$g_J=2$。饱和磁矩为
$g_JJ\mu_{\mathrm{B}}=5\mu_{\mathrm{B}}$。磁化率量的是
$\mu_{\mathrm{eff}}^{2}=g_J^{2}J(J+1)\mu_{\mathrm{B}}^{2}$，故
\begin{equation*}
g_J\sqrt{J(J+1)}\mu_{\mathrm{B}} = 2\times\sqrt{\frac{5}{2}\times\frac{7}{2}}\,\mu_{\mathrm{B}} = \sqrt{35}\,\mu_{\mathrm{B}} = 5.92\mu_{\mathrm{B}}.\tag{F.31}
\end{equation*}
饱和磁矩涉及 $\hat{J}_z$——它是沿特定方向的饱和磁矩测量；相比之下，磁化率涉及
$\hat{J}^{2}$。

\exer{4.1} 推导此结果有多种途径。其一从偶极子 1 的矢势出发：
\begin{equation*}
\mathbf{A}_1 = \frac{\mu_0}{4\pi}\nabla\times\frac{\boldsymbol{\mu}_1}{r},\tag{F.32}
\end{equation*}
它产生磁场
\begin{equation*}
\mathbf{B}_1 = \nabla\times\mathbf{A}_1 = \frac{\mu_0}{4\pi}\left[\nabla\left(\nabla\cdot\frac{\mu_1}{r}\right) - \nabla^{2}\left(\frac{\mu_1}{r}\right)\right],\tag{F.33}
\end{equation*}
给出能量
\begin{align*}
E &= -\boldsymbol{\mu}_2\cdot\mathbf{B}_1\\
&= -\frac{\mu_0}{4\pi}\left[(\boldsymbol{\mu}_1\cdot\nabla)(\boldsymbol{\mu}_1\cdot\nabla)\frac{1}{r} - (\boldsymbol{\mu}_1\cdot\boldsymbol{\mu}_2)\nabla^{2}\frac{1}{r}\right]\\
&= \frac{\mu_0}{4\pi r^{3}}\left[\boldsymbol{\mu}_1\cdot\boldsymbol{\mu}_2 - \frac{3(\boldsymbol{\mu}_1\cdot\mathbf{r})(\boldsymbol{\mu}_1\cdot\mathbf{r})}{r^{2}}\right].\tag{F.34}
\end{align*}
\footnote{原书如此。按推导，式 F.34 括号中第二、三项里的 $\boldsymbol{\mu}_1$
应为 $\boldsymbol{\mu}_2$（即 $3(\boldsymbol{\mu}_1\cdot\mathbf{r})(\boldsymbol{\mu}_2\cdot\mathbf{r})/r^{2}$）。——译注}
也可以用磁标势 $\phi_1=\mu_1\cdot r/4\pi r^{3}$、$B_1=-\mu_0\nabla\phi_1$ 推出
结果。

\exer{4.2} 1 \AA：（a）1.855 T（b）1.855 mT；10 \AA：（a）0.927 T（b）0.927 mT。

\exer{4.3} Fe（bcc）的晶格常数 $a=2.87$ \AA。最近邻距离为 $a/\sqrt{2}$，偶极--
偶极能量 $\sim(\mu_0\times(2.2\mu_{\text{B}})^{2}/4\pi r^{3})=30\ \mu$eV。
$\mathsf{J}\approx k_{\text{B}}T_{\text{C}}\sim0.09$ eV——大约大 3000 倍。

\exer{4.4} 带宽 $\sim0.05$ eV，库仑能 $\sim1$ eV。
$\mathsf{J}\sim-\frac{t^{2}}{U}=-2.5$ meV。故若 $|\mathsf{J}|=k_{\mathrm{B}}T_{\mathrm{N}}$，
则 $T_{\mathrm{N}}=29$ K。

\exer{4.5} 配分函数为
\begin{equation*}
Z = 1+\mathrm{e}^{-\beta\Delta}(\mathrm{e}^{\beta g\mu_{\mathrm{B}}B}+1+\mathrm{e}^{-\beta g\mu_{\mathrm{B}}B}).\tag{F.35}
\end{equation*}
由此可导出 F，进而 M：
\begin{equation*}
M = \frac{2n\beta(g\mu_{\mathrm{B}})^{2}B}{3+e^{\beta\Delta}},\tag{F.36}
\end{equation*}
$\chi\ll1$ 的结果随之而来。注意高温 $T\gg\Delta/k_{\mathrm{B}}$ 时
$\chi\to n\mu_0g^{2}\mu_{\mathrm{B}}^{2}/2k_{\mathrm{B}}T$——居里定律：所有能级
都被占据，磁化率是一个 $S=1$ 系统（$S(S+1)=2$）的 $\frac{3}{4}$ 与一个 $S=1$
系统（$S(S+1)=0$）\footnote{原书如此。由 $S(S+1)=0$ 可知应为 $S=0$——疑为原书
排印之误。——译注}的 $\frac{1}{4}$ 之和，即经典居里定律的 $\frac{3}{2}$。

低温 $T\ll\Delta/k_{\mathrm{B}}$ 时：$\Delta>0$ 则
$\chi\propto e^{-\Delta/k_{\mathrm{B}}T}\to0$（非磁——只有单态被占据）；$\Delta<0$
则 $\chi\to2n\mu_0g^{2}\mu_{\mathrm{B}}^{2}/3k_{\mathrm{B}}T$（对应被占据的三重态
$S(S+1)=2$，即经典居里定律的两倍）。

\exer{4.6} 提示：先证明
\begin{equation*}
\boldsymbol{\mu}_{I1}\cdot\boldsymbol{\mu}_{I2} \propto \hat{I}_{1z}\hat{I}_{2z} + \frac{1}{2}(\hat{I}_{1+}\hat{I}_{2-}+\hat{I}_{1-}\hat{I}_{2+})\tag{F.37}
\end{equation*}
并用 $\mathbf{r}=r(\sin\theta\cos\phi,\sin\theta\sin\phi,\cos\theta)$ 证明
\begin{equation*}
\boldsymbol{\mu}_I\cdot\hat{\mathbf{r}} \propto \cos\theta\hat{I}_z + \frac{1}{2}\sin\theta[e^{-i\phi}\hat{I}_{+}+e^{i\phi}\hat{I}_{-}].\tag{F.38}
\end{equation*}
结果随之很快得出。

\exer{4.8} 注意若 $\mathbf{S}=\mathbf{S}_1+\mathbf{S}_2+\mathbf{S}_3$，则
\begin{equation*}
\mathbf{S}^{2} = \mathbf{S}_1^{2}+\mathbf{S}_2^{2}+\mathbf{S}_3^{2} + 2(\mathbf{S}_1\cdot\mathbf{S}_2+\mathbf{S}_2\cdot\mathbf{S}_3+\mathbf{S}_3\cdot\mathbf{S}_1).\tag{F.39}
\end{equation*}
于是哈密顿量可以用 $\mathbf{S}^{2}$、$\mathbf{S}_1^{2}$、$\mathbf{S}_2^{2}$、
$\mathbf{S}_3^{2}$ 改写，给出 $E=\mathsf{J}(6-S(S+1))$。对 S 的各可能取值即得
所需能量。

\exer{5.1} $B_{\mathrm{mf}}=2.1\times10^{3}$ T。Fe 中单位体积原子数
$n=8.49\times10^{28}\,\mathrm{m^{-3}}$；用 $M=n\cdot2.2\mu_{\mathrm{B}}$ 得
$\mu_0M=2.2$ T。故 $B_{\mathrm{mf}}$ 比 $\mu_0M$ 大 $10^{3}$ 倍。

\exer{5.2} 磁化强度由
\begin{equation*}
\frac{M}{M_{\mathrm{s}}} = \mathrm{B}_S(y) = \alpha y - \beta y^{3}+\cdots,\tag{F.40}
\end{equation*}
给出，$\alpha$ 与 $\beta$ 是常数，而
\begin{equation*}
y = \frac{\gamma M}{TM_{\mathrm{s}}},\tag{F.41}
\end{equation*}
$\gamma$ 是由 $T_{\mathrm{C}}/\alpha$ 给出的正常数。$T_{\mathrm{C}}$ 附近这些方程
给出
\begin{equation*}
\left(\frac{M}{M_{\mathrm{s}}}\right)^{2} \approx \frac{\alpha^{3}}{\beta}\frac{(T_{\mathrm{C}}-T)}{T_{\mathrm{C}}}\tag{F.42}
\end{equation*}
用可由式 5.38 得到的 $\alpha$、$\beta$ 值即导出所需结果。

能量为 $E=-\frac{1}{2}\lambda M^{2}$；用式 F.42 可以证明在 $T_{\mathrm{C}}$ 处
\begin{equation*}
\frac{\mathrm{d}(M^{2})}{\mathrm{d}T} = -\frac{M_{\mathrm{s}}^{2}\alpha^{3}}{\beta T_{\mathrm{C}}},\tag{F.43}
\end{equation*}
由此可得比热的结果。

一般地容易证明
\begin{equation*}
C = \begin{cases} \frac{5}{2}nk_{\mathrm{B}}\left[\frac{(2S+1)^{2}-1}{(2S+1)^{2}+1}\right]\left(3\left(\frac{T}{T_{\mathrm{C}}}\right)^{2}-2\left(\frac{T}{T_{\mathrm{C}}}\right)\right) & T\leq T_{\mathrm{C}}\\ 0 & T>T_{\mathrm{C}}. \end{cases}\tag{F.44}
\end{equation*}

\exer{5.4} 能量为
\begin{equation*}
E = -2NS^{2}(\mathsf{J}_1\cos\theta + \mathsf{J}_2\cos2\theta),\tag{F.45}
\end{equation*}
能量极小 $\partial E/\partial\theta=0$ 给出 $\sin\theta=0$（从而 $\theta=0$ 或
$\pi$）或 $\cos\theta=-\mathsf{J}_1/4\mathsf{J}_2$。这三个解分别是铁磁（$\theta=0$）、
反铁磁（$\theta=\pi$）与螺旋磁（$\cos\theta=-\mathsf{J}_1/4\mathsf{J}_2$）；
代回能量方程即复现式 5.45。螺旋磁解只在 $|\mathsf{J}_1|<4|\mathsf{J}_2|$ 时可能。

哪个解能量最低可以通过比较各解的能量评估；或者通过计算 $\partial^{2}E/\partial\theta^{2}$
考察各解的稳定性。

\exer{5.5} $y\gg1$ 时
\begin{equation*}
\tanh y = \frac{e^{y}-e^{-y}}{e^{y}+e^{-y}} = \frac{1-e^{-2y}}{1+e^{-2y}} \approx 1-2e^{-2y}.\tag{F.46}
\end{equation*}
把它与
\begin{equation*}
\frac{M}{M_{\mathrm{s}}} = \tanh y,\tag{F.47}
\end{equation*}
以及 $T_{\mathrm{C}}$ 附近
\begin{equation*}
\frac{M}{M_{\mathrm{s}}} = \frac{T}{T_{\mathrm{C}}}y,\tag{F.48}
\end{equation*}
的事实联立，即可导出最终结果。

\exer{5.6} 回忆铁磁体的平均场理论：磁化强度由
\begin{equation*}
M = \frac{C}{\mu_0T}(\lambda M + \mu_0H),\tag{F.49}
\end{equation*}
给出，C 是居里常数。整理得
\begin{equation*}
M\left(1-\frac{\lambda C}{\mu_0T}\right) = \frac{HC}{T},\tag{F.50}
\end{equation*}
故磁化率为
\begin{equation*}
\chi = \frac{M}{H} = \frac{C}{T-\theta}\tag{F.51}
\end{equation*}
外斯温度 $\theta=\lambda C/\mu_0$。

对本题的铁磁体，每套子晶格的磁化强度为
\begin{equation*}
M_1 = \frac{C_1}{\mu_0T}(\mu_0H-\lambda M_2)\tag{F.52}
\end{equation*}
\begin{equation*}
M_2 = \frac{C_2}{\mu_0T}(\mu_0H-\lambda M_1)\tag{F.53}
\end{equation*}
即
\begin{equation*}
\begin{pmatrix} T & \lambda C_1/\mu_0\\ \lambda C_2/\mu_0 & T \end{pmatrix}\begin{pmatrix}M_1\\M_2\end{pmatrix} = H\begin{pmatrix}C_1\\C_2\end{pmatrix}.\tag{F.54}
\end{equation*}
此方程矩阵的行列式在 $T^{2}-\lambda^{2}C_1C_2/\mu_0=0$ 时为零，故可推断顺磁区为
$T>\theta$，$\theta=\lambda(C_1C_2)^{1/2}/\mu_0$。磁化率为
$\chi=(M_1+M_2)/H$；求逆式 F.54 得
\begin{equation*}
\begin{pmatrix}M_1\\M_2\end{pmatrix} = \frac{H}{T^{2}-\theta^{2}}\begin{pmatrix} T & -\lambda C_1/\mu_0\\ -\lambda C_2/\mu_0 & T \end{pmatrix}\begin{pmatrix}C_1\\C_2\end{pmatrix}\tag{F.55}
\end{equation*}
从而得到最终答案。

\exer{5.7} 第一部分直接由 $\hat{S}^{\pm}=\hat{S}^{x}\pm\mathrm{i}\hat{S}^{y}$ 的
应用得到。

对海森伯铁磁体，$\hat{S}_i^z\hat{S}_j^z$ 作用在基态上产生 $S^{2}$，基态是
$\hat{S}_i^z\hat{S}_j^z$ 的本征态。这样的项有 N 个，故 $E=-NS^{2}$。
$\hat{S}_i^{+}\hat{S}_j^{-}$ 与 $\hat{S}_i^{-}\hat{S}_j^{+}$ 不产生任何结果——
两种情形里升算符都湮灭基态。于是基态是二者的本征态、本征值为零。

对海森伯反铁磁体，$\hat{S}_i^{+}\hat{S}_j^{-}$ 与 $\hat{S}_i^{-}\hat{S}_j^{+}$
两项产生“想当然”基态之外的其他态，因此这一“想当然”的基态不是系统的本征态。

\exer{5.8} 可能的动量转移范围是 $0\,\mathrm{nm^{-1}}$ 到
$\sqrt{k_{\mathrm{in}}}=2\sqrt{2}\pi/\lambda=29.6\,\mathrm{nm^{-1}}$。

（a）100 K 时只看到化学（晶格）布拉格反射：
\begin{equation*}
Q = \left[\left(\frac{2\pi}{a}\right)^{2}(h^{2}+k^{2}) + \left(\frac{2\pi}{c}\right)^{2}l^{2}\right]^{1/2}.\tag{F.56}
\end{equation*}
对称性是 bcc：$h+k+l$ 为偶数时出现反射。

（b）10 K 时还看到磁布拉格反射。

两种情形观察到的反射列于下表。列出散射矢量 Q 的大小与入射束被散射的角 $2\theta$；
磁反射以星号标出，最后一列给出各反射的多重性 M。

\begin{table}
\centering
\begin{small}
\begin{tabular}{cccclcc}
\toprule
$h$ & $k$ & $l$ & Q（nm$^{-1}$） & $2\theta$（°） & & M\\
\midrule
0 & 1 & 0 & 12.566 & 34.915 & $*$ & 4\\
1 & 0 & 0 & 12.566 & 34.915 & $*$ & \\
1 & 1 & 0 & 17.772 & 50.208 & & 4\\
0 & 0 & 1 & 20.944 & 60.000 & $*$ & 2\\
0 & 1 & 1 & 24.425 & 71.337 & & 8\\
1 & 0 & 1 & 24.425 & 71.337 & & \\
0 & 2 & 0 & 25.133 & 73.740 & & 4\\
2 & 0 & 0 & 25.133 & 73.740 & & \\
1 & 1 & 1 & 27.468 & 81.952 & $*$ & 8\\
1 & 2 & 0 & 28.099 & 84.261 & $*$ & 8\\
2 & 1 & 0 & 28.099 & 84.261 & $*$ & \\
\bottomrule
\end{tabular}
\end{small}
\end{table}

\exer{5.9} （a）$\tilde{m}$ 的第一个表达式来自 $\tanh y$ 的定义。整理得
\begin{equation*}
x^{2} = \frac{1+\tilde{m}}{1-\tilde{m}},\tag{F.57}
\end{equation*}
取对数即式 5.57；用式 F.48 即得式 5.58。

（b）把式 5.60 整理成 x 的显式：
\begin{equation*}
(1-\tilde{m})x^{2} - \tilde{m}x - (1+\tilde{m}) = 0,\tag{F.58}
\end{equation*}
这是二次方程。$4-3\tilde{m}^{2}\geq0$ 时有实解，易得式 5.62。由
$0\leq\tilde{m}\leq1$ 知 $1\leq4-3\tilde{m}^{2}\leq2$，故 $x>0$ 需要式 5.62 中的
正根。于是
\begin{equation*}
y = \log(\tilde{m}+\sqrt{4-3\tilde{m}^{2}}) - \log(2(1-\tilde{m})),\tag{F.59}
\end{equation*}
且此情形 $y=3\tilde{m}T_{\mathrm{C}}/2T$，最终结果随之而来。

\exer{6.1} 平均能量为
\begin{equation*}
\langle E\rangle = -\frac{N\mathsf{J}}{4}\tanh\left(\frac{\mathsf{J}}{2k_{\mathrm{B}}T}\right),\tag{F.60}
\end{equation*}
每自旋热容为
\begin{equation*}
C = \frac{\mathsf{J}^{2}}{4k_{\mathrm{B}}T^{2}\cosh^{2}(\mathsf{J}/2k_{\mathrm{B}}T)}.\tag{F.61}
\end{equation*}
高温下它表现为 $\mathsf{J}^{2}/4k_{\mathrm{B}}T^{2}$，低温下如
$\mathsf{J}^{2}e^{-\mathsf{J}/k_{\mathrm{B}}T}/4k_{\mathrm{B}}T^{2}$。因此在
T=0 与 $T=\infty$ 处都为零、中间有一个宽极大——没有尖峰，这印证了纯一维模型没有
相变的论点。

\exer{6.2}
\begin{equation*}
\text{(b)}\quad \hbar\omega(\mathbf{q}) = S[\mathsf{J}(\mathbf{0})-\mathsf{J}(\mathbf{q})+K(\mathbf{0})].
\end{equation*}
\begin{equation*}
\text{(c)(i)}\quad \hbar\omega(\mathbf{q}) = 2S\mathsf{J}_0(1-\cos qa) + 2SK_0.
\end{equation*}
\begin{equation*}
\text{(ii)}\quad \hbar\omega(\mathbf{q}) = 2S\mathsf{J}_0(2-\cos q_xa-\cos q_ya) + 4SK_0.
\end{equation*}
\begin{equation*}
\text{(iii)}\quad \hbar\omega(\mathbf{q}) = 8S\mathsf{J}_0\left(1-\cos\frac{q_xa}{2}\cos\frac{q_ya}{2}\cos\frac{q_za}{2}\right) + 8SK_0.
\end{equation*}
\exer{6.3} 小 q 时用 $\cos qa\approx1-(qa)^{2}/2$ 得
（i）$\hbar\omega = S(\mathsf{J}_0a^{2}q^{2}+2K_0)$
（ii）$\hbar\omega = S(\mathsf{J}_0a^{2}q^{2}+4K_0)$
（iii）$\hbar\omega = S(\mathsf{J}_0a^{2}q^{2}+8K_0)$

伊辛情形（$\mathsf{J}_0=0$）：$\hbar\omega=2dSK_0\equiv\Delta$，d 对（i）（ii）
（iii）分别为 1、2、4。此模型中自旋波无论波矢为何都“花费”同样的能量 $\Delta$。
自旋波数目正比于 $e^{-\Delta/k_{\mathrm{B}}T}$，其能量正比于
$\Delta e^{-\Delta/k_{\mathrm{B}}T}$。比热容易算得正比于
\begin{equation*}
\frac{\Delta^{2}}{k_{\mathrm{B}}T^{2}}\,\mathrm{e}^{-\Delta/k_{\mathrm{B}}T},
\end{equation*}
与习题 6.1 一致（本题的 $\Delta$ 即该题的 $\mathsf{J}$）。

海森伯情形（$K_0=0$）：$\hbar\omega=Dq^{2}$，自旋刚度 $D=S\mathsf{J}_0a^{2}$。
自旋波数目 $N_{\mathrm{s}}$ 为
\begin{equation*}
N_{\mathrm{s}} \propto \int\frac{\mathrm{d}^{n}q}{\mathrm{e}^{Dq^{2}/k_{\mathrm{B}}T}-1},\tag{F.62}
\end{equation*}
小 q 时
\begin{equation*}
\frac{1}{\mathrm{e}^{Dq^{2}/k_{\mathrm{B}}T}-1} \approx \frac{k_{\mathrm{B}}T}{Dq^{2}}.\tag{F.63}
\end{equation*}
积分在一、二维发散（$\mathrm{d}^{n}q$ 分别为 $\mathrm{d}q$ 与 $2\pi q\,\mathrm{d}q$），
三维收敛（$\mathrm{d}^{n}q=4\pi q^{2}\mathrm{d}q$）。因此纯（各向同性）海森伯
模型在一、二维不可能有长程序。

\exer{6.4} 反铁磁情形可以循铁磁情形的推导，但须记住自己位于哪套子晶格。得到
\begin{equation*}
\begin{aligned}
\hbar\dot{S}_j^{+} &= 2\mathrm{i}\mathsf{J}S[2S_j^{+}+S_{j-1}^{+}+S_{j+1}^{+}] \quad j \text{ 为奇}\\
\hbar\dot{S}_j^{+} &= -2\mathrm{i}\mathsf{J}S[2S_j^{+}+S_{j-1}^{+}+S_{j+1}^{+}] \quad j \text{ 为偶}.
\end{aligned}\tag{F.64}
\end{equation*}
代入
\begin{equation*}
S_j^{+} = \begin{cases} u\,\mathrm{e}^{\mathrm{i}(qja-\omega t)} & j \text{ 为奇}\\ v\,\mathrm{e}^{\mathrm{i}(qja-\omega t)} & j \text{ 为偶}, \end{cases}\tag{F.65}
\end{equation*}
得
\begin{equation*}
\mathrm{i}\hbar\omega\begin{pmatrix}u\\v\end{pmatrix} = 4\mathrm{i}\mathsf{J}S\begin{pmatrix} 1 & \cos qa\\ -\cos qa & -1 \end{pmatrix}\begin{pmatrix}u\\v\end{pmatrix},\tag{F.66}
\end{equation*}
从而
\begin{equation*}
\begin{vmatrix} 1-\frac{\hbar\omega}{4\mathsf{J}S} & \cos qa\\ \cos qa & 1+\frac{\hbar\omega}{4\mathsf{J}S} \end{vmatrix} = 0\tag{F.67}
\end{equation*}
故 $\hbar\omega=4\mathsf{J}S\sin qa$。

\exer{6.5} 用 $\partial F/\partial M=0$，少许整理即得
\begin{equation*}
M^{2} = -\frac{2a_0(T-T_{\mathrm{C}})}{4b} + \frac{\mu_0H}{4bM},\tag{F.68}
\end{equation*}
正是所需形式。把 $M^{2}$ 对 H/M 作图得直线，$M^{2}$ 轴上的截距在 $T=T_{\mathrm{C}}$
处变号。

\exer{6.6} 合适单位下模型为
\begin{equation*}
\epsilon = \frac{1}{2}\sin^{2}(\theta-\phi) - h\cos\phi.\tag{F.69}
\end{equation*}
$\theta=0$ 时能量极小给出
\begin{equation*}
\sin\phi(\cos\phi+h) = 0\tag{F.70}
\end{equation*}
解为 $\phi=0, \pi, \cos^{-1}(-h)$。计算 $\partial^{2}E/\partial\theta^{2}$\footnote{原书
如此。按题意能量 $\epsilon$ 为 $\phi$ 的函数，此处应为 $\partial^{2}E/\partial\phi^{2}$
——疑为原书排印之误。——译注}可考察
这些解的稳定性。$\theta=\pi/2$ 时能量极小给出
\begin{equation*}
\sin\phi(h-\cos\phi) = 0\tag{F.71}
\end{equation*}
解为 $\phi=0, \pi, \cos^{-1}(h)$。

\exer{6.7} 磁化强度的分量为
\begin{equation*}
M_x = |\mathbf{M}|\sin\frac{\pi r}{R}\tag{F.72}
\end{equation*}
\begin{equation*}
M_y = |\mathbf{M}|\cos\frac{\pi r}{R},\tag{F.73}
\end{equation*}
用
\begin{equation*}
\frac{\mathrm{d}r}{\mathrm{d}x} = \frac{\mathrm{d}}{\mathrm{d}x}\sqrt{x^{2}+y^{2}+z^{2}} = \frac{x}{r}\tag{F.74}
\end{equation*}
之类的结果可以求出 $M_x$ 与 $M_y$ 的所有偏导数。于是
\begin{equation*}
(\nabla M_x)^{2}+(\nabla M_y)^{2} = |\mathbf{M}|^{2}\frac{\pi^{2}}{R^{2}},\tag{F.75}
\end{equation*}
最终结果由此而来。

\exer{6.8} 用合理的单位，25 Gbits\,in$^{-2}$ 即 $3.9\times10^{13}$ bits\,m$^{-2}$。

（a）取 $\pi a_0^{2}$ 为一个氢原子的面积，得 $3.4\times10^{-7}$ bits（氢原子
面积）$^{-1}$。

（b）莎士比亚的剧本多在两万到四万词之间；每词平均约 5 个字母，共约 40 部戏，
大约 4 Mbytes 即 32 Mbits（每字节 8 比特）。拿一张英国邮票算，我把每张邮票算出
约 600 份全集——用数据压缩还能翻倍。“我对这些数目字全不在行”（《哈姆雷特》
第二幕第二场）。

DNA 每碱基对存 1 比特，约每立方纳米 1 比特。若做成一张薄片，相当于约
$6\times10^{5}$ Gbits\,in$^{-2}$；但 DNA 的巨大优势在于它能卷曲，是体积存储技术。

\exer{6.9} 最好的做法是把磁场取在 $-z$ 方向——不失一般性。（为何 $-z$？这样
基态所有自旋平行于 z。电子带负电、自旋与其磁矩反平行，因而将与外场反平行排列。）
你会发现每个磁振子添加能量 $+g\mu_{\mathrm{B}}B$。

\exer{6.11} 方程
\begin{equation*}
\langle q|i\rangle = \frac{1}{\sqrt{N}}\sum_j e^{-i\mathbf{q}\cdot\mathbf{R}_j}\langle j|i\rangle\tag{F.76}
\end{equation*}
可以用 $\langle j|i\rangle=\delta_{ij}$ 简化，结果随之而来。

下一部分的提示：$S^{+}$ 会湮灭它所作用的态——除非它作用的自旋小于 S。

本题表明每个自旋都有一个垂直于磁化方向的小横向分量；这些横向分量随空间的变化
正如图 6.13 所画。

\exer{6.12} 需要求解
\begin{equation*}
\Delta E = 0 = \int_{-\infty}^{\infty}\left[2A\frac{\partial\theta}{\partial z}\frac{\partial}{\partial z}(\delta\theta) + \frac{\partial f(\theta)}{\partial\theta}\delta\theta\right]\mathrm{d}z\tag{F.77}
\end{equation*}
积分的第一部分分部积分后给出
\begin{equation*}
\frac{\partial f}{\partial\theta} - 2A\frac{\partial^{2}f}{\partial z^{2}} = 0.\tag{F.78}
\end{equation*}
题目本身细心地引导读者完成其余部分。

\exer{6.13} 由结构中壁的长度带来的壁能量密度约为 $\sigma_{\mathrm{W}}L/D$。
小三角形内部的自旋耗费各向异性能量。每个三角形面积 $D^{2}/4$，每个水平长度 D
中有两个，总代价为 $2\times K\times(D^{2}/4)\times(1/D)=KD/2$。把这些项相加并
求导即得所需结果。$D=20\ \mu$m。

\exer{7.1} 用 $g(E_{\mathrm{F}})=3n/2E_{\mathrm{F}}$，第一个结果很快得到。要考虑
温度依赖，值得回顾关于费米函数的一些数学细节。设
$h(E)=\int_0^E g(E')\mathrm{d}E'$（$g(E)=\mathrm{d}h/\mathrm{d}E$），考虑积分
\begin{align*}
I &= \int_{0}^{\infty}\frac{\mathrm{d}h}{\mathrm{d}E}f(E)\,\mathrm{d}E\\
&= [f(E)g(E)]_{0}^{\infty} - \int_{0}^{\infty}h(E)\frac{\mathrm{d}f}{\mathrm{d}E}\,\mathrm{d}E\\
&= -\int_{0}^{\infty}h(E)\frac{\mathrm{d}f}{\mathrm{d}E}\,\mathrm{d}E.\tag{F.79}
\end{align*}
令 $x=(E-\mu)/k_{\mathrm{B}}T$，则
\begin{equation*}
\frac{\mathrm{d}f}{\mathrm{d}E} = -\frac{1}{k_{\mathrm{B}}T}\frac{e^{x}}{(e^{x}+1)^{2}}.\tag{F.80}
\end{equation*}
写
\begin{equation*}
h(E) = \sum_{s=0}^{\infty}\frac{x^{s}}{s!}\left(\frac{\mathrm{d}^{s}h}{\mathrm{d}x^{s}}\right)_{x=0},\tag{F.81}
\end{equation*}
有
\begin{equation*}
I = \sum_{s=0}^{\infty}\frac{1}{s!}\left(\frac{\mathrm{d}^{s}h}{\mathrm{d}x^{s}}\right)_{x=0}\int_{-E_{\mathrm{F}}/k_{\mathrm{B}}T}^{\infty}\frac{x^{s}e^{x}\,\mathrm{d}x}{(e^{x}+1)^{2}}.\tag{F.82}
\end{equation*}
把下限换成 $-\infty$ 可简化其中的积分。奇 s 时积分为零；偶 s 时
\begin{align*}
\int_{-\infty}^{\infty}\frac{x^{s}\mathrm{e}^{x}\,\mathrm{d}x}{(\mathrm{e}^{x}+1)^{2}} &= 2\int_{0}^{\infty}\frac{x^{s}\mathrm{e}^{x}\,\mathrm{d}x}{(\mathrm{e}^{x}+1)^{2}}\\
&= 2\int_{0}^{\infty}\mathrm{d}x\sum_{n=0}^{\infty}\mathrm{e}^{x}x^{s}\left[(n+1)(-1)^{n+1}\mathrm{e}^{-nx}\right]\\
&= 2\sum_{n=1}^{\infty}(-1)^{n+1}n\int_{0}^{\infty}x^{s}\mathrm{e}^{-nx}\,\mathrm{d}x\\
&= 2(s!)\sum_{n=1}^{\infty}\frac{(-1)^{n+1}}{n^{s}}\\
&= 2(s!)(1-2^{1-s})\zeta(s),\tag{F.83}
\end{align*}
$\zeta(s)$ 是黎曼 zeta 函数（$\zeta(0)=-\frac{1}{2}$、$\zeta(2)=\frac{\pi^{2}}{6}$、
$\zeta(4)=\frac{\pi^{4}}{90}$）。

于是积分为
\begin{align*}
I &= \sum_{s=0,\,s\text{ 偶}}^{\infty}2\left(\frac{\mathrm{d}^{s}h}{\mathrm{d}x^{s}}\right)_{x=0}(1-2^{1-s})\zeta(s)\\
&= h + \frac{\pi^{2}}{6}\left(\frac{\mathrm{d}^{2}h}{\mathrm{d}x^{2}}\right)_{x=0} + \frac{7\pi^{4}}{360}\left(\frac{\mathrm{d}^{4}h}{\mathrm{d}x^{4}}\right)_{x=0} + \cdots\\
&= \int_{-\infty}^{\mu}g(E)\,\mathrm{d}E + \frac{\pi^{2}}{6}(k_{\mathrm{B}}T)^{2}\left(\frac{\mathrm{d}g}{\mathrm{d}E}\right)_{E=\mu}\\
&\quad + \frac{7\pi^{4}}{360}(k_{\mathrm{B}}T)^{4}\left(\frac{\mathrm{d}^{3}g}{\mathrm{d}E^{3}}\right)_{E=\mu} + \cdots.\tag{F.84}
\end{align*}
这意味着数密度为
\begin{equation*}
n = \int_{0}^{\mu}g(E)\,\mathrm{d}E + \frac{\pi^{2}}{6}(k_{\mathrm{B}}T)^{2}\left(\frac{\mathrm{d}g}{\mathrm{d}E}\right)_{E=\mu} + \cdots.\tag{F.85}
\end{equation*}
一阶近似下
\begin{equation*}
\int_{0}^{\mu}g(E)\,\mathrm{d}E = \int_{0}^{E_{\mathrm{F}}}g(E)\,\mathrm{d}E + (\mu-E_{\mathrm{F}})g(E_{\mathrm{F}}),\tag{F.86}
\end{equation*}
故对恒定的电子数，n 在一切温度取同值，从而
\begin{equation*}
\mu = E_{\mathrm{F}} - \frac{\pi^{2}}{6}(k_{\mathrm{B}}T)^{2}\frac{g'(E_{\mathrm{F}})}{g(E_{\mathrm{F}})}.\tag{F.87}
\end{equation*}
磁场的作用是把能级移动 $\pm\mu_{\mathrm{B}}B$。每个自旋态的数密度为
\begin{align*}
n_{\pm} &= \frac{1}{2}\int_{0}^{E_{\mathrm{F}}}g(E)\,\mathrm{d}E + (\mu-E_{\mathrm{F}}\mp\mu_{\mathrm{B}}B)g(E_{\mathrm{F}}\mp\mu_{\mathrm{B}}B)\\
&\quad + \frac{\pi^{2}}{6}(k_{\mathrm{B}}T)^{2}\left(\frac{\mathrm{d}g}{\mathrm{d}E}\right)_{E=\mu\mp\mu_{\mathrm{B}}B} + \cdots\tag{F.88}
\end{align*}
用近似关系
\begin{equation*}
g(E_{\mathrm{F}}\mp\mu_{\mathrm{B}}B) = g(E_{\mathrm{F}}) \mp \mu_{\mathrm{B}}B\left(\frac{\mathrm{d}g}{\mathrm{d}E}\right)_{E=\mu}\tag{F.89}
\end{equation*}
与
\begin{equation*}
\left(\frac{\mathrm{d}g}{\mathrm{d}E}\right)_{E=\mu\mp\mu_{\mathrm{B}}B} = \left(\frac{\mathrm{d}g}{\mathrm{d}E}\right)_{E=\mu} \mp \mu_{\mathrm{B}}B\left(\frac{\mathrm{d}^{2}g}{\mathrm{d}E^{2}}\right)_{E=\mu}\tag{F.90}
\end{equation*}
得
\begin{equation*}
n_{+}-n_{-} = -\frac{\pi^{2}}{6}(k_{\mathrm{B}}T)^{2}\mu_{\mathrm{B}}B\left[\left(\frac{g'}{g}\right)^{2}-\frac{g''}{g}\right] + \mu_{\mathrm{B}}Bg(E_{\mathrm{F}}).\tag{F.91}
\end{equation*}
故磁化率为
\begin{equation*}
\chi = \mu_0\mu_{\mathrm{B}}g(E_{\mathrm{F}})\left[1-\frac{\pi^{2}}{6}(k_{\mathrm{B}}T)^{2}\left[\left(\frac{g'}{g}\right)^{2}-\frac{g''}{g}\right]\right].\tag{F.92}
\end{equation*}
$g(E)\propto E^{1/2}$ 时易证
\begin{equation*}
\frac{g'}{g} = \frac{1}{2E_{\mathrm{F}}},\qquad \left(\frac{g'}{g}\right)^{2}-\frac{g''}{g} = \frac{1}{2E_{\mathrm{F}}^{2}}\tag{F.93}
\end{equation*}
（在费米能处取值）。最终结果随之而来。对 $E_{\mathrm{F}}\sim10$ eV 的金属
（如 Sn），$T_{\mathrm{F}}\sim10^{5}$ K，修正约为 $10^{-5}$。

\exer{7.2} （a）这在式 7.23 的推导中已经示明。

（b）交叉温度为 $2T_{\mathrm{F}}/3=27$ K。泡利顺磁磁化率为 $1.3\times10^{-8}$。
朗道抗磁磁化率为
\begin{equation*}
\chi_{\mathrm{L}} = -\left(\frac{m_{\mathrm{e}}}{m^{*}}\right)^{2}\frac{\chi_{\mathrm{P}}}{3} = -4.3\times10^{-7}.\tag{F.94}
\end{equation*}

\exer{7.3} 用 $A_{\phi}=\frac{1}{2}Br$，薛定谔方程成为
\begin{align*}
\Bigg[\frac{-\hbar^{2}}{2m_{\mathrm{e}}}\left(\frac{\partial^{2}}{\partial r^{2}}+\frac{1}{r}\frac{\partial}{\partial r}+\frac{1}{r^{2}}\frac{\partial^{2}}{\partial\phi^{2}}+\frac{\partial^{2}}{\partial z^{2}}\right) & - \frac{\mathrm{i}\hbar eB}{2m_{\mathrm{e}}}\frac{\partial}{\partial\phi}\\
& + \left.\frac{e^{2}B^{2}r^{2}}{8m_{\mathrm{e}}}\right]\psi(r,\phi,z) = E\psi(r,\phi,z).\tag{F.95}
\end{align*}
寻找如下形式的解
\begin{equation*}
\psi(r,\phi,z) = R(r)e^{\mathrm{i}l\phi}e^{\mathrm{i}k_z z}\tag{F.96}
\end{equation*}
导致
\begin{align*}
\Bigg(-\frac{\hbar^{2}}{2m_{\mathrm{e}}}\left(\frac{\partial^{2}}{\partial r^{2}}+\frac{1}{r}\frac{\partial}{\partial r}\right) - \frac{\hbar^{2}l^{2}}{2m_{\mathrm{e}}} + \frac{1}{2}m_{\mathrm{e}}\omega_{\mathrm{c}}^{2}r^{2} + \frac{\hbar k_z^{2}}{2m_{\mathrm{e}}}\Bigg)R(r) &= \left(E-\frac{l}{2}\hbar\omega_{\mathrm{c}}\right)R(r).\tag{F.97}
\end{align*}
解为
\begin{equation*}
\psi(r,\theta,z) = \exp\left[\mathrm{i}(l\phi+k_z z)-\frac{eBr^{2}}{4\hbar}\right]r^{|l|}L_{n-1}^{|l|}\left(\frac{eBr^{2}}{2\hbar}\right),\tag{F.98}
\end{equation*}
$L_{\alpha}^{\beta}(x)$ 是连带拉盖尔多项式（附录 C.5），能量为
\begin{equation*}
E = \left(n+\frac{l}{2}+\left|\frac{l}{2}\right|-\frac{1}{2}\right)\hbar\omega_{\mathrm{c}} + \frac{\hbar k_z^{2}}{2m_{\mathrm{e}}}.\tag{F.99}
\end{equation*}
在这一所谓对称规范中，所有态都绕空间中同一点旋转——以牺牲平移对称换取了旋转
对称的显现。

\exer{7.5} 提示：
\begin{equation*}
\frac{z+1}{z-1} = 1+\frac{2}{z-1}\tag{F.100}
\end{equation*}
由此可以证明
\begin{equation*}
1+\frac{1-z^{2}}{2z}\log\left(\frac{z+1}{z-1}\right) = \frac{2}{3z^{2}}+O\left(\frac{1}{z^{3}}\right).\tag{F.101}
\end{equation*}
另外 $\log(-1)=\mathrm{i}\pi$，故
\begin{equation*}
\log\left|\frac{x+1}{x-1}\right| - \log\left(\frac{x+1}{x-1}\right) = \mathrm{i}\pi.\tag{F.102}
\end{equation*}

\exer{7.6} （a）这在式 7.23 的推导中已经示明。（见习题 7.2(a)。）

（b）从下式出发很有用：
\begin{equation*}
\chi_{\mathbf{q}} = \mu_0\mu_{\mathrm{B}}^{2}\frac{2m_{\mathrm{e}}}{\hbar^{2}}I
\end{equation*}
I 是积分
\begin{equation*}
I = \int_{|\mathbf{k}|<k_{\mathrm{F}}}g(\mathbf{k})\,\mathrm{d}^{n}k\left[\frac{1}{(k+q)^{2}-k^{2}} + \frac{1}{(k-q)^{2}-k^{2}}\right].\tag{F.103}
\end{equation*}
一维中
\begin{align*}
\int_{|\mathbf{k}|<k_{\mathrm{F}}}\frac{\mathrm{d}k}{(k+q)^{2}-k^{2}} &= \int_{-k_{\mathrm{F}}}^{k_{\mathrm{F}}}\frac{\mathrm{d}k}{q(q+2k)}\\
&= \frac{1}{2q}\int_{q-2k_{\mathrm{F}}}^{q+2k_{\mathrm{F}}}\frac{\mathrm{d}u}{u}\\
&= \frac{1}{2q}\int_{|q-2k_{\mathrm{F}}|}^{q+2k_{\mathrm{F}}}\frac{\mathrm{d}u}{u}\\
&= \frac{1}{2q}\log\left|\frac{q+2k_{\mathrm{F}}}{q-2k_{\mathrm{F}}}\right|.\tag{F.104}
\end{align*}
结果随之而来。

二维中要考虑的积分是
\begin{align*}
I &= \int_{0}^{k_{\mathrm{F}}}k\,\mathrm{d}k\int_{0}^{2\pi}\frac{\mathrm{d}\theta}{q(q-2k\cos\theta)} + \frac{\mathrm{d}\theta}{q(q+2k\cos\theta)}\\
&= 2\int_{0}^{k_{\mathrm{F}}}k\,\mathrm{d}k\int_{0}^{2\pi}\frac{\mathrm{d}\theta}{q^{2}-4k^{2}\cos\theta}\tag{F.105}
\end{align*}
而
\begin{equation*}
\int_{0}^{2\pi}\frac{\mathrm{d}\theta}{q^{2}-4k^{2}\cos\theta} = \begin{cases} \frac{2\pi}{q(q^{2}-4k^{2})^{1/2}} & q>2k_{\mathrm{F}}\\ 0 & q<2k_{\mathrm{F}} \end{cases}\tag{F.106}
\end{equation*}
于是积分成为
\begin{equation*}
I = \begin{cases} \pi\left[1-\sqrt{1-\left(\frac{2k_{\mathrm{F}}}{q}\right)}\right] & q>2k_{\mathrm{F}}\\ \pi & q<2k_{\mathrm{F}} \end{cases}\tag{F.107}
\end{equation*}
最终结果随之而来。
